\scanpage{26}
un morphisme, on va prolonger $u$ \deleted{de} façon \uline{canonique} en un \unsure{hom.} $\bar u$ de $X$ dans $Y$ compatible avec les augmentations. Si on a déjà défini $\bar u^n:X^n\unsure{-|\to}Y^n$ \deleted{\unsure{et}} pour un $n$ $(n\geq-1)$, montrons comment continuer : \unsure{on a} $s^{n+1}\in\mathrm{Hom}(\unsure{\overline{X^{n+1}}|\widetilde{X^{n+1}}},\overline{X^n})$,

\diagram{\begin{tikzcd}[ampersand replacement=\&]
X^n\arrow[r,"d^n"]\arrow[d,"\bar u^n"']\& X^{n+1}\\
Y^n\arrow[r,dash,"d^n"]\&Y^{n+1}\\
\unsure{\vert|1}\&
\end{tikzcd}}

\deleted{il définit} \inserted{soit} $s'^{n+1}\simeq\mathrm{Hom}(\deleted{\unsure{\widetilde X}}X^{n+1},T(X^n))$ définit un \deleted{\unsure{hom}} morphisme $\deleted{\unsure{t}}T(\bar u^n)\circ s'^{n+1}:X^{n+1}-T(Y^n)$, d'où, \deleted{\illegible{1}} \inserted{composant avec le} morphisme $t^n:T(Y^n)-Y^{n+1}$, un morphisme \deleted{$\unsure{X^i}$} $t^nT(\bar u^n)s'^{n+1}:X^{n+1}\to Y^{n+1}$, posons

\[
\bar u^{n+1}=t^nT(\bar u^n)s'^{n+1}
\]

\deleted{On a $\bar u^{n+1}d^n=t^nT(\bar u^n)s'^{n+1}d^n=-t^nT(\bar u^n)d$}

On vérifie que ça marche \unsure{\ldots}.

\uline{Résolution canonique} $C^*(A)$, définie par
\[
C^{n+1}(A)=T(C^n(A)/\mathrm{Im}\,\widehat C^{n-1}(A))\quad (\deleted{\unsure{en|\text{avec}}}\ \text{posant }\widehat C^{-1}(A)=A)
\]

\deleted{$\unsure{C^*(A)|C^\bullet}$} C'est à la fois une $s$ et une $t$-résolution, c'est un foncteur en $A$. L'hom fonctoriel $C^*(A)-C^*(B)$ associé à un hom $A-B$ est en fait celui défini grâce à la $s$-structure de $C^*(A)$ et la $t$-structure de $C^*(B)$.

% paper watermark: ARGENTINE TYPEWRITER (mirrored, vertical)
% archive sheet number: 25
